\documentclass[11pt,a4paper]{article}
\usepackage[margin=25mm]{geometry}
\usepackage[T1]{fontenc}
\usepackage{lmodern,microtype,amsmath,amssymb,amsthm,mathtools,booktabs,xcolor}
\usepackage[hidelinks,pdftitle={How the zeros of deformed monotone Hurwitz polynomials approach their limits},pdfauthor={Anonymous}]{hyperref}
\newtheorem{theorem}{Theorem}[section]
\newtheorem{proposition}[theorem]{Proposition}
\newtheorem{corollary}[theorem]{Corollary}
\newtheorem{lemma}[theorem]{Lemma}
\newtheorem{conjecture}[theorem]{Question}
\theoremstyle{definition}\newtheorem{definition}[theorem]{Definition}
\theoremstyle{remark}\newtheorem{remark}[theorem]{Remark}\newtheorem{example}[theorem]{Example}
\DeclareMathOperator{\sgn}{sgn}
\numberwithin{equation}{section}
\title{How the zeros of deformed monotone Hurwitz polynomials approach their limits}
\author{Anonymous\\[3pt]{\small DOI \href{https://doi.org/10.5281/zenodo.24546686}{10.5281/zenodo.24546686}}}
\date{November 2026}
\begin{document}
\maketitle
\begin{abstract}
Coulter, Do and Moskovsky (CDM) introduced polynomials $\vec H^t_{g,n}(\mu)$ that count monotone factorisations of a permutation of cycle type $\mu$ into transpositions, with a weight $t$ for each distinct larger letter used. They conjectured that the zeros are real, that they interlace, and (Conjecture~3.15) that as the genus $g$ grows each zero converges to an explicit limit, from below if the limit is less than $-1$ and from above otherwise. The limits follow from the leading large-genus asymptotics of Accadia, Lewa\'nski and Ruzza. We compute the next term for every profile, and for any number of profiles. It depends on the profile only through its number of fixed points, as in the analogous result of Yang for simple Hurwitz numbers. This term fixes the side and rate of approach of every zero, so the eventual form of Conjecture~3.15 holds, with an explicit genus threshold. Real-rootedness and interlacing at all sufficiently large genera (the large-genus cases of Conjectures~3.13 and~3.14) already follow from the leading term, and we record them explicitly. Monotonicity cannot start at genus zero: exact examples in degrees $4$, $6$, $8$ and $10$ move away from their limits. For every profile of degree at most $10$ we certify real-rootedness at every genus, and we determine exactly the genera at which each zero moves towards its limit. Further results give the third term (fixed points and $2$-cycles), show that each early term depends only on short cycles, identify the limiting joint law of the larger letters used, and explain the exact coincidences in CDM's numerical tables.
\end{abstract}

\tableofcontents

\section{Introduction}\label{sec:intro}
This section states the problem in plain terms, says why it is of interest, separates what was known from what is new, and gives a map of the paper.

\paragraph{The problem.}
Write a permutation $\sigma$ of $\{1,\dots,d\}$ as a product of transpositions, $\sigma=(a_1\,b_1)\cdots(a_r\,b_r)$ with $a_i<b_i$ for each $i$. Call the factorisation \emph{monotone} if the larger letters never decrease, $b_1\le b_2\le\cdots\le b_r$, and \emph{transitive} if the transpositions together connect all $d$ letters. Its \emph{hive number} is the number of distinct larger letters $b_i$. Counting transitive monotone factorisations of a fixed $\sigma$ with weight $t^{\text{hive number}}$ gives a polynomial in $t$. A transitive factorisation of $\sigma$ into $r$ transpositions corresponds to a branched cover of the sphere, whose genus $g$ is given by $r=2g-2+d+n$, where $n$ is the number of cycles of $\sigma$; large genus therefore means long factorisations. The cycle type $\mu$ of $\sigma$ is called its (ramification) \emph{profile}.

Coulter, Do and Moskovsky \cite{cdm} met these counts when integrating over the complex Grassmannian of $M$-dimensional subspaces of $\mathbb C^N$: the same weighted counts without the transitivity condition are the coefficients of the large-$N$ expansion of its Weingarten function at fixed $t=1-M/N$. The transitive counts, divided by the product of the cycle lengths, are their connected versions, which CDM call deformed monotone Hurwitz numbers. At $t=1$ they are the classical monotone Hurwitz numbers \cite{cdm}. The zeros studied here lie at $t<0$, that is $M>N$, so the root conjectures are structural statements about the polynomials, not about a physical Grassmannian.

CDM proved that the coefficients are symmetric and unimodal. Their data pointed to much more. Their Conjecture~3.13 says the polynomials have only real zeros. Conjecture~3.14 says the zeros interlace when one cycle grows by one. Conjecture~3.15 describes the zeros for a fixed cycle type $\mu$ of size $d$ as the genus $g$ grows. The zeros should tend to
\[
 -\frac{d-2}{1},\ -\frac{d-3}{2},\ \dots,\ -\frac{1}{d-2},\ 0,
\]
and the convergence should be increasing from below for limits less than $-1$ and decreasing from above for nonzero limits greater than $-1$.

\paragraph{An example.}
For the $4$-cycle the polynomial (before CDM's normalisation) is $t+3t^2+t^3$ at genus $0$, $15t+40t^2+15t^3$ at genus $1$ and $161t+413t^2+161t^3$ at genus~$2$. Its zeros other than $0$ are $-2.618,-0.382$, then $-2.215,-0.451$, then $-2.086,-0.479$. They approach the predicted limits $-2$ and $-\tfrac12$, the first from below and the second from above. CDM's tables at genus $10$ to $20$ in degree $7$ show the same behaviour.

Why a definite side? For long factorisations the count is dominated by the trivial representation of $S_d$, whose largest Jucys--Murphy eigenvalue $d-1$ gives the limit. The first correction, with the next eigenvalue $d-2$, comes from the trivial and standard representations and from the factorisations that never move one fixed point of $\sigma$. It depends only on the number of fixed points, and its sign at each limit decides the side.

\paragraph{Why it matters.}
Real zeros, interlacing and monotone root motion are the signature of a hidden positive structure: a polynomial with nonnegative coefficients and only real zeros is, after scaling, the probability generating function of a sum of independent Bernoulli variables. Proving such statements usually needs a mechanism. CDM proved real-rootedness and interlacing for one-cycle profiles in genera $0$ and $1$, using a recursion in the degree of a special three-term shape, and showed that no recursion of that shape exists in genus~$2$ \cite[Proposition~3.17 and Remark~3.18]{cdm}. The large-genus regime is a natural place to look for a different one.

\paragraph{What was known.}
Accadia, Lewa\'nski and Ruzza (ALR) \cite{alr} determined the leading large-genus asymptotics of a wide class of weighted Hurwitz numbers. Their theorem covers counts weighted by a rational function $G$. CDM's polynomials are the case $G(z)=(1+(t-1)z)/(1-z)$, and the theorem then says that, for long factorisations, $h_r$ is approximately $\frac{2}{d!}(d-1)^rA_{d-1}(t)$, where $A_{d-1}$ is the explicit polynomial \eqref{eq:A} below. Its zeros are exactly CDM's limits, and they are distinct. Hence the limits in Conjecture~3.15, and real-rootedness and interlacing at all sufficiently large genera, already follow from ALR. The leading term does not say from which side the zeros approach. At $t=1$, Yang \cite[Theorem~1.1 and Proposition~1.2]{yang} proved that monotone Hurwitz numbers are exponential polynomials in the genus and identified the leading term with an error of order $(d-2)^r$. For \emph{simple} Hurwitz numbers, Do, He and Robertson conjectured that the coefficients between the two largest exponentials vanish \cite[Conjecture~4.1]{dhr}; Yang proved this and showed that the next coefficient is $-d$ times the number of fixed points \cite[Theorem~1.3]{yang}.

\paragraph{What is new.}
Fix $d$ and $\mu$, let $f$ be the number of fixed points of $\mu$, and put $D=d-1$.
\begin{enumerate}
\item \emph{The second term.} Theorem~\ref{th:two} gives the next term of the large-genus expansion exactly, as a polynomial in $t$: it is $-(D+f-1)(1+t)A_{D-1}(t)$ times $(D-1)^r$, where $r$ is the length of the factorisation. The theorem covers any finite list of prescribed profiles.
\item \emph{Direction of approach.} Theorem~\ref{th:roots} shows that each moving zero differs from its limit by an explicit multiple of $((D-1)/D)^r$, with the sign CDM predicted. So every zero eventually moves strictly towards its limit from the predicted side. Theorem~\ref{th:uniform} makes ``eventually'' explicit: $r\ge200d^2\log(20d)$ suffices.
\item \emph{Every genus in small degree.} Monotonicity cannot start at genus $0$ in general (Theorem~\ref{th:reverse}). For all $132$ profiles of degrees $4$ to $10$, Theorem~\ref{th:census} proves real-rootedness at every genus and lists exactly when each zero moves towards its limit. Only eight profiles ever move away, all in even degree, and never after genus~$2$.
\item \emph{Higher terms and short cycles.} The coefficient of $(D-j)^r$ depends only on the numbers of cycles of length at most $j$, for $j<d/2$ (Theorem~\ref{th:cycles}). The third term is explicit in terms of fixed points and $2$-cycles (Theorem~\ref{th:third}).
\item \emph{A probabilistic reading.} As the genus grows, the set of larger letters used by a uniformly random factorisation converges to independent Bernoulli choices, letter $b$ being used with probability $(b-1)/(d-1)$ (Proposition~\ref{pr:occupancy}).
\item \emph{CDM's tables explained.} An identity read off from CDM's cut-and-join recursion explains why some rows of CDM's Figure~2 coincide exactly (Proposition~\ref{pr:fixed}), and the second and third terms reproduce every row (Table~\ref{tab:fig2}).
\end{enumerate}
The method is standard: the content-product formula for Jucys--Murphy elements, inclusion--exclusion over set partitions for transitivity, and partial fractions. The new input is the explicit evaluation of the second and third poles, and the finite certificates.

\paragraph{Roadmap.}
Section~\ref{sec:setup} fixes notation and proves three lemmas on the generating function. Section~\ref{sec:two} proves the two-term expansion. Section~\ref{sec:zeros} derives the behaviour of the zeros at large genus, with an explicit threshold, and compares with CDM's tables. Section~\ref{sec:small} treats every genus in small degree. Section~\ref{sec:higher} gives the higher terms. Section~\ref{sec:prob} gives the probabilistic reading. Section~\ref{sec:open} lists open questions, and Appendix~\ref{app:code} describes the reproducible computations.

\section{Set-up and three lemmas on the generating function}\label{sec:setup}
This section defines the polynomials, the leading polynomial $A_k$, and the generating function, and proves the three facts about its poles that every later section uses.

\subsection{Definitions and notation}
Fix $d\ge1$ and $\sigma\in S_d$ with cycle type $\mu=(\mu_1,\dots,\mu_n)$. For $r\ge0$ let
\[
 h_r(\mu;t)=\sum t^{\,\#\{b_1,\dots,b_r\}},
\]
summed over the transitive monotone factorisations $(a_1\,b_1)\cdots(a_r\,b_r)=\sigma$ of length $r$. The sum depends only on $\mu$ \cite{cdm}. The Riemann--Hurwitz relation links length and genus: $r=2g-2+d+n$. In CDM's normalisation \cite[Definition~3.1]{cdm}
\[
 \vec H^t_{g,n}(\mu_1,\dots,\mu_n)=\frac{h_r(\mu;t)}{\mu_1\cdots\mu_n},
\]
which has the same zeros as $h_r$. Throughout we write
\[
 D=d-1,\qquad q=\frac{D-1}{D},\qquad f=m_1(\mu),\qquad c=m_2(\mu),
\]
where $m_a(\mu)$ is the number of cycles of length $a$. CDM proved that $h_r$ has degree exactly $D$ in $t$, has no constant term, and satisfies $h_r(t)=t^d h_r(1/t)$ \cite[Proposition~3.9]{cdm}. In particular $t=0$ is always a zero, and $t=-1$ is a zero whenever $D$ is even.

For $k\ge1$ define the \emph{leading polynomials}
\begin{equation}\label{eq:A}
 A_k(t)=t\prod_{j=1}^{k-1}\frac{k-j+jt}{k-j}=\prod_{j=1}^{k}\Bigl(t-1+\frac kj\Bigr).
\end{equation}
The two forms agree because $\prod_{j=1}^{k-1}j/(k-j)=1$. Thus $A_k$ is monic with the $k$ distinct zeros $x_{k,j}=1-k/j$, $1\le j\le k$. For example $A_1=t$, $A_2=t(t+1)$ and $A_3=t(t+2)(t+\tfrac12)$. The zeros of $A_D$ are CDM's limits. The nonzero zeros of $A_{k-1}$ strictly interlace those of $A_k$, because
\begin{equation}\label{eq:interlace}
 1-\frac kj<1-\frac{k-1}{j}<1-\frac k{j+1}\qquad(1\le j<k-1).
\end{equation}

\paragraph{Several profiles.}
Sections~\ref{sec:two} and~\ref{sec:higher} allow $N\ge1$ prescribed cycle types $\mu^{(1)},\dots,\mu^{(N)}$ of $S_d$. Choose $\sigma_i$ independently and uniformly from these classes, and count monotone words $\tau_1\cdots\tau_r$ with
\[
 \sigma_1\cdots\sigma_N\,\tau_1\cdots\tau_r=1
\]
for which the group generated by all the $\sigma_i$ and $\tau_j$ is transitive, with weight $t$ to the hive number. Let $T_r(t)$ be the expectation of this weighted count. For $N=1$, $T_r=h_r(\mu;t)$, because $\sigma^{-1}$ has the same cycle type as $\sigma$. The admissible lengths have $(-1)^r=\prod_i\sgn\sigma_i$, and the genus is given by $r=2g-2+(2-N)d+\sum_i\ell(\mu^{(i)})$. Write $f_i$ for the number of fixed points of $\mu^{(i)}$ and $m_{i,a}$ for its number of $a$-cycles.

\subsection{Three lemmas}
Let $F_c(z)=\dfrac{1+(t-1)cz}{1-cz}$, and for a Young diagram $\lambda$ write $c(\square)$ for the content (column minus row) of a box. Let $\chi^\lambda$ be the irreducible character and $\dim\lambda$ its degree.

\begin{lemma}[Content product]\label{lem:content}
Let $U(z;t)=\sum_r U_r(t)z^r$ be defined like $\sum_rT_r(t)z^r$ but without the transitivity condition. Then
\begin{equation}\label{eq:char}
 U(z;t)=\frac{1}{d!}\sum_{\lambda\vdash d}(\dim\lambda)^2\prod_{i=1}^N\frac{\chi^\lambda(\sigma_i)}{\dim\lambda}\prod_{\square\in\lambda}F_{c(\square)}(z).
\end{equation}
\end{lemma}
\begin{proof}
Let $J_b=\sum_{a<b}(a\,b)$ be the Jucys--Murphy elements of the group algebra. The words whose larger letters all equal $b$ contribute $1+t\sum_{m\ge1}(zJ_b)^m=F_{J_b}(z)$, since a nonempty run earns one factor $t$. A monotone word is a sequence of such runs in increasing $b$, so the generating series is $\prod_bF_{J_b}(z)$. This product is a symmetric function of the $J_b$ and therefore acts on the irreducible module $V_\lambda$ by the scalar $\prod_{\square\in\lambda}F_{c(\square)}(z)$. Each class average acts on $V_\lambda$ by $\chi^\lambda(\sigma_i)/\dim\lambda$. The coefficient of the identity of a central element is its trace in the regular representation divided by $d!$, which gives \eqref{eq:char}. For $N=1$ this is \cite[Proposition~3.3]{cdm}, and the general product appears in \cite[Eq.~(1.26)]{alr} with another normalisation.
\end{proof}

\begin{lemma}[Inclusion--exclusion over set partitions]\label{lem:mobius}
For a set partition $\pi$ of $\{1,\dots,d\}$, let $U_\pi$ be the generating function of tuples in which every $\sigma_i$ and every transposition preserves each block of $\pi$. Then
\begin{equation}\label{eq:mobius}
 \sum_rT_r(t)z^r=\sum_{\pi}(-1)^{|\pi|-1}(|\pi|-1)!\;U_\pi(z;t),
\end{equation}
and $U_\pi$ is the probability that all $\sigma_i$ preserve the blocks of $\pi$ times a finite weighted sum of products, over the blocks $P\in\pi$, of the functions \eqref{eq:char} for $S_{|B|}$ with the restricted cycle types.
\end{lemma}
\begin{proof}
A tuple preserves the blocks of $\pi$ exactly when its orbit partition $\rho$ refines $\pi$. Möbius inversion in the lattice of set partitions gives $\sum_{\pi\ge\rho}(-1)^{|\pi|-1}(|\pi|-1)!=[\rho\text{ has one block}]$, which proves \eqref{eq:mobius}. Fix $\pi$ and a choice of cycle types of the restrictions of the $\sigma_i$ to each block. Conditional on preserving $\pi$ with these restricted types, the restrictions are independent and uniform in their block classes, because conjugation by $\prod_{P\in\pi}S_P$ preserves the event and acts transitively on each choice. A monotone word that preserves $\pi$ is the merge of one monotone word in each block. Since distinct blocks have disjoint sets of larger letters, the merge is unique and the hive numbers add. After order-preserving relabelling of each block, the block factor is the function \eqref{eq:char} for that block.
\end{proof}

\begin{lemma}[Poles]\label{lem:poles}
The function \eqref{eq:mobius} is rational in $z$. Its poles lie at $z=\pm1/k$ with $1\le k\le D$, and the pole at $\pm1/k$ has order at most $\lfloor d/(k+1)\rfloor$. Its polynomial part is constant. Under $z\mapsto-z$ it changes by the factor $\prod_i\sgn\sigma_i$, so on admissible lengths the two poles $\pm1/k$ contribute equally. Consequently, on admissible lengths $r\ge1$, each coefficient of $T_r$ is a sum over $1\le k\le D$ of $k^r$ times a polynomial in $r$ of degree less than $\lfloor d/(k+1)\rfloor$. If $2(k+1)>d$, a term of \eqref{eq:mobius} has a pole at $1/k$ only if one block has at least $k+1>d/2$ letters, so this block is unique, and the pole is simple.
\end{lemma}
\begin{proof}
Each $F_c$ is a Möbius function of $z$ with a single pole at $1/c$ ($c\ne0$) and is bounded at infinity, so every product in \eqref{eq:char} is rational with constant polynomial part. A box of content $k\ge1$ sits in a row of length at least $k+1$, and a row contains at most one box of each content. Hence a diagram with $s$ boxes of content $k$ has at least $s$ rows of length at least $k+1$, and so at least $s(k+1)$ boxes. In a product over the blocks of $\pi$ the block sizes sum to $d$, so the total number of boxes of content $k$ is at most $\lfloor d/(k+1)\rfloor$. The same holds for content $-k$ by transposing diagrams. Each transposition changes the sign of a permutation, so a word of length $r$ multiplies it by $(-1)^r$, and the generating function is even or odd in $z$ according to $\prod_i\sgn\sigma_i$. A rational function whose poles have these orders has coefficients of the stated form for $r\ge1$. For the last assertion, a box of content $k$ needs a block of size at least $k+1$, two such blocks would need $2(k+1)>d$ letters, and two boxes of content $k$ in one block would need $2(k+1)>d$ boxes.
\end{proof}

\section{The second term of the large-genus expansion}\label{sec:two}
This section proves the two-term expansion for any finite list of profiles, and states the single-profile case. Define
\begin{equation}\label{eq:LK}
 L=D^{2-N}\prod_{i=1}^N(f_i-1),\qquad K=d^{2-N}\prod_{i=1}^Nf_i,\qquad B(t)=1+L-K-\Bigl(D+\frac LD\Bigr)t .
\end{equation}

\begin{theorem}[Two-term expansion]\label{th:two}
Fix $d\ge3$ and any finite list of profiles. As $r\to\infty$ through admissible lengths,
\begin{equation}\label{eq:general}
 \frac{d!}{2}\,T_r(t)=D^rA_D(t)+(D-1)^rA_{D-1}(t)B(t)+E_r(t),
\end{equation}
where for $d\ge4$ every coefficient of $E_r$ is $O_{d,\boldsymbol\mu}\bigl((1+r)^{d}(D-2)^r\bigr)$, and for $d=3$ we have $E_r=0$ for every admissible $r\ge1$.
\end{theorem}

\begin{corollary}[One profile]\label{cor:single}
For one profile with $f$ fixed points, $B(t)=-(D+f-1)(1+t)$, so
\begin{equation}\label{eq:single}
 \frac{d!}{2}\,h_r(\mu;t)=D^rA_D(t)-(D+f-1)(1+t)(D-1)^rA_{D-1}(t)+E_r(t).
\end{equation}
\end{corollary}
\begin{proof}
With $N=1$, $L=D(f-1)$ and $K=df$, so $B(t)=1+D(f-1)-df-(D+f-1)t=-(D+f-1)(1+t)$.
\end{proof}

\begin{proof}[Proof of Theorem~\ref{th:two}]
By Lemma~\ref{lem:poles} it suffices to find the coefficients of $(1-Dz)^{-1}$ and $(1-(D-1)z)^{-1}$ in the partial fraction expansion of \eqref{eq:mobius}; the poles at $-1/D$ and $-1/(D-1)$ double these on admissible lengths, which accounts for the factor $2$.

\emph{Step 1: the pole at $1/D$.} Content $D$ occurs only in the one-row diagram $(d)$ of the one-block partition. This diagram has $\chi^{(d)}=1$, dimension $1$ and contents $0,1,\dots,D$. Multiplying by $1-Dz$ and setting $z=1/D$ gives $\frac1{d!}\prod_{c=1}^{D-1}F_c(1/D)$ times the numerator $1+(t-1)=t$ of $F_D$. Since $F_c(1/D)=(D-c+ct)/(D-c)$, this is $A_D(t)/d!$.

\emph{Step 2: the pole at $1/(D-1)$ in the one-block term.} Content $D-1$ occurs in the diagrams $(d)$ and $(d-1,1)$. Both contain the contents $0,1,\dots,D-1$ once each, which contribute $A_{D-1}(t)$ as in Step~1. The remaining box has content $D$ in $(d)$ and $-1$ in $(d-1,1)$. The representation $(d-1,1)$ is the permutation representation minus the trivial one, so it has dimension $D$ and character $f_i-1$ at $\sigma_i$. Its weight in \eqref{eq:char}, relative to the row, is $D^2\prod_i\frac{f_i-1}{D}=L$. The one-block contribution to the coefficient is therefore $\frac{A_{D-1}(t)}{d!}\bigl[F_D(\tfrac1{D-1})+LF_{-1}(\tfrac1{D-1})\bigr]=\frac{A_{D-1}(t)}{d!}\bigl[(1-Dt)+L(1-t/D)\bigr]$.

\emph{Step 3: proper set partitions.} A proper partition reaches content $D-1=d-2$ only if it has a block of size $d-1$, so it is a singleton and its complement. There are $d$ of these, each with Möbius coefficient $-1$. A given letter is fixed by all $\sigma_i$ with probability $\prod_if_i/d$. Conditional on that, the restrictions to the complement are uniform in their classes, and by Step~1 in degree $d-1$ the complement contributes $A_{D-1}(t)/(d-1)!$ at this pole. The total is $-d\prod_i(f_i/d)\,A_{D-1}(t)/(d-1)!=-KA_{D-1}(t)/d!$.

Adding Steps~2 and~3 gives the coefficient $A_{D-1}(t)B(t)/d!$. By Lemma~\ref{lem:poles} all remaining poles have $|z|\ge1/(D-2)$ when $d\ge4$, with order at most $d$, which gives the remainder. For $d=3$ the only poles are $\pm1$ and $\pm\frac12$, and both are simple, so the expansion is exact for $r\ge1$.
\end{proof}

\begin{remark}[What the second term means]\label{rem:mechanism}
The Jucys--Murphy elements act by contents, so the exponential rates in the genus are the contents of Young diagrams. The largest content, $D$, comes only from the trivial representation; that is the leading term. The next content, $D-1$, has three sources: the trivial representation again, the standard representation $(d-1,1)$, whose character counts fixed points, and the correction for words that never move one fixed point of $\sigma$, which are not transitive and must be subtracted. All three see the profile only through fixed points. This is the monotone, $t$-deformed counterpart of Yang's formula for simple Hurwitz numbers, where the first subleading coefficient is $-d$ times the number of fixed points \cite[Theorem~1.3]{yang}; Yang's theorem also proves the vanishing conjectured in \cite[Conjecture~4.1]{dhr}.
\end{remark}

\begin{remark}[Relation to the leading term]\label{rem:alr}
Step~1 is the specialisation of ALR's leading asymptotics \cite[Lemma~2.3, Theorem~6 and Theorem~C]{alr} to $G(z)=(1+(t-1)z)/(1-z)$; the connected leading term needs no correction because no proper set partition reaches content $D$. ALR note that further terms can be obtained by subtracting principal parts \cite[Remark~2.2]{alr}. Steps~2 and~3 carry this out at the second pole.
\end{remark}

\section{The zeros at large genus}\label{sec:zeros}
This section turns the two-term expansion into statements about zeros. It proves the eventual form of CDM's Conjecture~3.15 with an explicit threshold, records the consequences for Conjectures~3.13 and~3.14, and compares the predictions with CDM's tables.

\subsection{Direction of approach}
For $1\le j<D$ write $x_j=x_{D,j}=1-D/j$ and
\[
 R_j=\frac{A_{D-1}(x_j)}{A_D'(x_j)} .
\]
By \eqref{eq:interlace}, $A_{D-1}(x_j)$ and $A_D'(x_j)$ both have the sign $(-1)^{D-j}$, so $R_j>0$.

\begin{theorem}[Zeros of one profile]\label{th:roots}
Let $d\ge3$ and fix $\mu$. For all sufficiently large genera, $h_r(\mu;t)$ has $D$ simple real zeros. Labelled in increasing order, the $j$th zero $\alpha_{g,j}$ ($1\le j<D$) satisfies
\begin{equation}\label{eq:shift}
 \alpha_{g,j}=x_j+(D+f-1)(1+x_j)R_j\,q^{r}+o(q^{r}),\qquad r=2g-2+d+n .
\end{equation}
Hence a zero with limit $x_j<-1$ eventually increases strictly to $x_j$ from below, and a zero with limit $-1<x_j<0$ eventually decreases strictly to $x_j$ from above. The zero $0$ is exact at every genus, and so is $-1$ when $D$ is even. Moreover, at every sufficiently large common genus, every profile of degree $d$ strictly interlaces every profile of degree $d+1$ after the common zero $0$ is removed.
\end{theorem}

The limits, eventual real-rootedness and eventual interlacing need only the leading term and so already follow from ALR \cite{alr}. The direction and rate of approach need the second term. In CDM's terms, Theorem~\ref{th:roots} proves the eventual form of Conjecture~3.15 and the large-genus cases of Conjectures~3.13 and~3.14; the interlacing statement covers the pairs $\mu$ and $\mu+e_i$ of Conjecture~3.14 and in fact any pair of adjacent degrees. In degree $d=2$ the polynomial is $t$ at every positive admissible length, and in degree $3$ it is $\frac{2^r-f-1}{3}\,t(1+t)$ by \eqref{eq:single}, so there is no moving zero.

Theorem~\ref{th:roots} is the case $N=1$ of the following statement.

\begin{corollary}[Zeros for several profiles]\label{cor:multi}
For all sufficiently large admissible $r$, $T_r$ has $D$ simple real zeros, one of which is $0$. The zero $\alpha_{r,j}$ near $x_j$ satisfies
\[
 \alpha_{r,j}=x_j-B(x_j)R_j\,q^r+o(q^r).
\]
If $B(x_j)\ne0$, it is eventually strictly monotone in steps of two in $r$, approaching from below if $B(x_j)>0$ and from above if $B(x_j)<0$. Every list of profiles in degree $d$ eventually interlaces every list in degree $d+1$ at the same genus, strictly after division by $t$.
\end{corollary}
\begin{proof}
Divide \eqref{eq:general} by $D^r$; the result converges coefficientwise to $A_D$. Choose disjoint intervals about the $D$ simple zeros of $A_D$ with opposite signs of $A_D$ at the two ends of each. Eventually the scaled polynomial has the same signs there, so it has a zero in each interval. Its degree is at most $D$, so these are all its zeros, and each is simple. Every word of positive length uses a larger letter, so $0$ is a zero. At a nonzero limit $x_j$, evaluate \eqref{eq:general} at the nearby zero $\alpha$ and divide by the difference quotient $(A_D(\alpha)-A_D(x_j))/(\alpha-x_j)$, which tends to $A_D'(x_j)\ne0$. This gives the displayed displacement. Writing $\kappa=-B(x_j)R_j$, we get $\alpha_{r+2,j}-\alpha_{r,j}=\kappa(q^2-1)q^r+o(q^r)$, which proves eventual strict monotonicity. Finally, the strict inequalities $1-\frac{D+1}{j}<1-\frac Dj<1-\frac{D+1}{j+1}$ ($1\le j<D$) between the nonzero zeros of $A_D$ and $A_{D+1}$ persist under small perturbation, also when the two lengths at a common genus differ by a constant.
\end{proof}

\begin{proof}[Proof of Theorem~\ref{th:roots}]
Apply Corollary~\ref{cor:multi} with $N=1$ and $B(t)=-(D+f-1)(1+t)$. The displacement is $(D+f-1)(1+x_j)R_jq^r$, which is negative for $x_j<-1$ and positive for $-1<x_j<0$. When $D$ is even, $d$ is odd, and $h_r(t)=t^dh_r(1/t)$ gives $h_r(-1)=0$.
\end{proof}

\begin{remark}[Several profiles]\label{rem:double}
For two fixed-point-free profiles, $L=1$ and $K=0$, so $B(t)=2-(D+1/D)t>0$ on the negative axis: every negative zero, including those in $(-1,0)$, eventually approaches its limit from below. This is a different pattern from the single-profile case. It is related to the ``double'' enumeration that CDM mention in \cite[Remark~3.6]{cdm}. The case $B(x_j)=0$, which Corollary~\ref{cor:multi} leaves open, does occur. For $d=7$ and the profiles $(6,1)$ and $(3,1,1,1,1)$ we have $L=0$, $K=4$ and $B(t)=-3-6t$, which vanishes at the limit $x_4=-\frac12$. The direction there is decided by a later term.
\end{remark}

\subsection{An explicit threshold}
The next theorem gives a length beyond which all the conclusions of Theorem~\ref{th:roots} hold, uniformly over profiles of degree $d$. Write
\[
 S(t)=(D+f-1)(1+t)A_{D-1}(t),
\]
so that \eqref{eq:single} reads $\frac{d!}2h_r=D^rA_D-(D-1)^rS+E_r$.

\begin{theorem}[Uniform threshold]\label{th:uniform}
Let $d\ge3$. For every profile of degree $d$ and every admissible length
\[
 r\ge200\,d^2\log(20d)
\]
(natural logarithm), $h_r(\mu;t)$ has $D$ simple real zeros, and from each such admissible length $r$ to the next one, $r+2$, each moving zero moves strictly towards its limit from the side given in Theorem~\ref{th:roots}. More precisely, with $b_j=S(x_j)/A_D'(x_j)$, the $j$th zero lies in the interval with endpoints $x_j+b_jq^r(1\pm\frac1{4d})$.
\end{theorem}
\begin{proof}
Put $M=(20d)^{10d}$, $\theta=(D-\frac32)/(D-1)$ and $F_r=d!\,h_r/(2D^r)$. We first show that, on the admissible parity and uniformly for complex $|t|\le2d$,
\begin{equation}\label{eq:uniform-error}
 \bigl|F_r(t)-A_D(t)+q^rS(t)\bigr|\le M(q\theta)^r .
\end{equation}
Take the circle $|z|=1/(D-\frac32)$ in \eqref{eq:mobius}. For each integer content $|c|\le D$ we have $|cz|\le4$ and $|1-cz|\ge1/(2d)$, because $D-\frac32$ is a half-integer. Hence $|F_c(z)|\le(1+4(2d+1))\,2d\le20d^2$. Since $|\chi^\lambda|\le\dim\lambda$ and $\sum_\lambda(\dim\lambda)^2=b!$ in $S_b$, every block factor in Lemma~\ref{lem:mobius} has modulus at most $(20d^2)^b$. There are at most $d^d$ set partitions, and $d!$ and each Möbius coefficient are at most $d^d$. So $\frac{d!}2$ times \eqref{eq:mobius} has modulus at most $(20d^5)^d$ on the circle. The principal parts at $\pm1/D$ and $\pm1/(D-1)$ have numerators $\pm\frac12A_D$ and $\pm\frac12S$ in the scaling of \eqref{eq:single}, and their denominators have modulus at least $1/(2d)$ on the circle. As $A_k$ is monic with $k$ zeros in $[-d,0]$, $|A_D|\le(3d)^d$ and $|S|\le2d(1+2d)(3d)^d$ for $|t|\le2d$. Subtracting the four principal parts leaves a function analytic on the closed disc with boundary modulus less than $M$. Cauchy's estimate, division by $D^r$ and the parity pairing give \eqref{eq:uniform-error}.

Fix a moving zero $x=x_j\ne-1$ and set $b=S(x)/A_D'(x)$. Distinct zeros of $A_D$ and $A_{D-1}$ are rationals with denominators at most $d$, so they are at least $d^{-2}$ apart, and \eqref{eq:interlace} excludes a common nonzero zero. Also $|1+x|\ge1/d$. The product formulas give
\[
 |A_D'(x)|\ge d^{-2(D-1)},\qquad |S(x)|\ge d^{-2D+1}>1/M,\qquad |b|\le2d(1+2d)(3d)^dd^{2(D-1)}<M .
\]
Differentiating the products gives $|A_D''|,|S'|\le M$ on $|t|\le2d$. The length bound gives $q^r,\theta^r\le e^{-r/(2d)}\le M^{-10}$. In particular the two test points $u_\pm=x+bq^r(1\pm\varepsilon)$, $\varepsilon=1/(4d)$, lie within $2M^{-9}$ of $x$. By Taylor's theorem and \eqref{eq:uniform-error},
\[
 \bigl|F_r(u_\pm)\mp\varepsilon S(x)q^r\bigr|\le4M^3q^{2r}+M(q\theta)^r<\varepsilon|S(x)|q^r,
\]
where the last inequality follows from $4M^4q^r+M^2\theta^r\le4M^{-6}+M^{-8}<\varepsilon$. The opposite signs give a zero between $u_-$ and $u_+$. The intervals for different $j$ are disjoint by the separation bound. The exact zeros $0$ and, for even $D$, $-1$ account for the remaining indices, and $F_r(1)>0$ by \eqref{eq:uniform-error} since $A_D(1)\ge1$. So $F_r$ is a nonzero polynomial of degree at most $D$ with $D$ distinct real zeros, all simple. Finally $q^2(1+\varepsilon)<1-\varepsilon$ for $d\ge3$, so the intervals at lengths $r$ and $r+2$ are disjoint and ordered towards $x$.
\end{proof}

The bound is far from optimal. For $d=7$ it asks for $r\ge48{,}429$, while the exact computations of Section~\ref{sec:small} show monotone approach from genus $0$ for every profile of degree $7$.

\subsection{An exact identity and CDM's tables}
CDM's Figure~2 lists the zeros at genus $20$ for all fifteen profiles of degree $7$. Several rows agree to all nine printed decimals. The next identity explains the exact coincidences.

\begin{proposition}[Removing a fixed point]\label{pr:fixed}
For every genus $g$ and every nonempty $\nu=(\nu_1,\dots,\nu_n)$,
\[
 \vec H^t_{g,n+1}(1,\nu_1,\dots,\nu_n)=\sum_{i=1}^n(\nu_i+1)\,\vec H^t_{g,n}(\nu_1,\dots,\nu_i+1,\dots,\nu_n).
\]
In particular, if all parts of $\nu$ equal $k$, then $(1,k^{n})$ and $(k+1,k^{n-1})$ have proportional polynomials at every genus, and so the same zeros.
\end{proposition}
\begin{proof}
Put $\mu_1=1$ in CDM's cut-and-join recursion \cite[Theorem~3.7]{cdm}. The term with factor $(\mu_1-1)$ vanishes, and the sum over $\alpha+\beta=\mu_1$ with $\alpha,\beta\ge1$ is empty. Only the join terms remain, which is the identity. Combinatorially: since $h_r$ depends only on the cycle type, we may take the fixed point to be the letter $d$; the last transposition $(a\,d)$ of a transitive word then always joins it to another cycle.
\end{proof}

This gives the four exact coincidences in CDM's Figure~2: $(7)$ with $(6,1)$, $(4,3)$ with $(3,3,1)$, $(3,2,2)$ with $(2,2,2,1)$, and $(2,1^5)$ with $(1^7)$. Other rows agree to first order. At a fixed genus the displacement \eqref{eq:shift} is proportional to $(D+f-1)q^{n}$. For example $(5,2)$, $(4,3)$, $(4,2,1)$ and $(3,3,1)$ all give $5q^2$, and they differ only in the ninth decimal place, where the third term of Theorem~\ref{th:third} separates them. Table~\ref{tab:fig2} shows how well the expansions predict all fifteen rows.

\begin{table}[ht]
\centering\small
\begin{tabular}{lcccc}
\toprule
profile & $f$ & $\alpha_{20,1}$ (exact, 12 digits) & two-term error & three-term error\\
\midrule
$(7)$, $(6,1)$ & $0,1$ & $-5.00126794186$ & $1.7\times10^{-4}$ & $5.5\times10^{-10}$\\
$(5,2)$ & $0$ & $-5.00105643527$ & $1.4\times10^{-4}$ & $3.3\times10^{-10}$\\
$(5,1,1)$ & $2$ & $-5.00123268479$ & $1.6\times10^{-4}$ & $5.2\times10^{-10}$\\
$(4,3)$, $(3,3,1)$ & $0,1$ & $-5.00105643832$ & $1.3\times10^{-4}$ & $6.6\times10^{-11}$\\
$(4,2,1)$ & $1$ & $-5.00105643629$ & $1.3\times10^{-4}$ & $2.4\times10^{-10}$\\
$(4,1,1,1)$ & $3$ & $-5.00117392864$ & $1.5\times10^{-4}$ & $4.4\times10^{-10}$\\
$(3,2,2)$, $(2,2,2,1)$ & $0,1$ & $-5.00088023534$ & $1.1\times10^{-4}$ & $4.0\times10^{-11}$\\
$(3,2,1,1)$ & $2$ & $-5.00102706599$ & $1.3\times10^{-4}$ & $1.5\times10^{-10}$\\
$(3,1^4)$ & $4$ & $-5.00110049213$ & $1.4\times10^{-4}$ & $3.0\times10^{-10}$\\
$(2,2,1,1,1)$ & $3$ & $-5.00097811783$ & $1.2\times10^{-4}$ & $1.2\times10^{-10}$\\
$(2,1^5)$, $(1^7)$ & $5,7$ & $-5.00101890606$ & $1.3\times10^{-4}$ & $1.9\times10^{-10}$\\
\bottomrule
\end{tabular}
\caption{Degree $7$, genus $20$: the smallest zero of each profile (CDM's Figure~2, first column) and the relative error in its displacement from the limit $-5$ when the zero of the truncated two-term and three-term expansions is used instead. Rows joined by a comma are equal by Proposition~\ref{pr:fixed}. The other three moving zeros behave in the same way; the worst relative errors over all sixty entries are $1.7\times10^{-4}$ and $5.5\times10^{-10}$. All values agree with CDM's nine printed (truncated) decimals.}\label{tab:fig2}
\end{table}

\section{Every genus in small degree}\label{sec:small}
This section asks whether the conclusions of Section~\ref{sec:zeros} hold from genus $0$. They do not in general, but in degrees up to $10$ they fail only in a few explicit cases, which we determine exactly.

\subsection{Zeros that move away}
\begin{theorem}[Zeros moving away from their limits]\label{th:reverse}
For $\mu=(2,2)$ the two nonzero zeros equal their limits $-2,-\frac12$ at genus $0$ and move away at genus $1$. For $\mu=(2,2,2)$ a zero already strictly below its limit $-\frac32$ at genus $0$ moves strictly further below at genus $1$.
\end{theorem}
\begin{proof}
For $\sigma=(1\,2)(3\,4)$ the characters of $S_4$ on the shapes $(4),(3,1),(2,2),(2,1,1),(1^4)$ are $1,-1,2,-1,1$ and their dimensions are $1,3,2,3,1$. An intransitive word factors independently over the two $2$-cycles, so its generating series is $t^2z^2/(1-z^2)^2$. Subtracting this from \eqref{eq:char} and putting $u=z^2$ gives
\[
 \frac{2tu^2\bigl[(2-5u+3u^2)(1+t^2)+(5-11u-6u^2)t\bigr]}{(1-u)^2(1-4u)(1-9u)} .
\]
The coefficients of $u^2$ and $u^3$ are $4t+10t^2+4t^3$ and $50t+128t^2+50t^3$. At genus $0$ the nonzero zeros are $-2$ and $-\frac12$. At genus $1$ they are $(-32\pm\sqrt{399})/25$, outside $[-2,-\frac12]$.

In degree $6$ the polynomials of type $(2,2,2)$ at lengths $7$ and $9$ are $8t(7t^4+53t^3+96t^2+53t+7)$ and $24t(112t^4+797t^3+1434t^2+797t+112)$; they are obtained by the direct word count described in Appendix~\ref{app:code}. Divide the quartics by $t^2$ and put $y=t+1/t$. They become $7y^2+53y+82$ and $112y^2+797y+1210$. Both are negative at $y=-3$ and positive at $y=-13/6$, and at $y=-24/11$ their values are $-38/121$ and $514/121$. So their larger roots lie on opposite sides of $-24/11$, with the genus-$1$ root smaller. As $t\mapsto t+1/t$ increases on $(-\infty,-1)$ and maps $-\frac32$ to $-\frac{13}6$, both corresponding zeros lie below $-\frac32$, and the genus-$1$ zero is further from it.
\end{proof}

So Conjecture~3.15 is true in its eventual form but fails if monotonicity is required from genus $0$. In the $(2,2)$ case the genus-$0$ zeros sit exactly at their limits, so the failure is degenerate; the $(2,2,2)$ case is a genuine reversal.

\subsection{A census of degrees four to ten}
\begin{theorem}[All genera, degree at most ten]\label{th:census}
Let $\mu$ be any of the $132$ profiles of degree $4\le d\le10$. Then:
\begin{enumerate}
\item at every genus $g\ge0$, $h_r(\mu;t)$ has $D$ simple real zeros, so CDM's Conjecture~3.13 holds for $\mu$;
\item every moving zero below $-1$ increases strictly towards its limit, and every moving zero in $(-1,0)$ decreases strictly towards its limit, at every step from genus $g$ to $g+1$, with the following exceptions. For $(2,2)$, $(2,2,2)$, $(4,2,2)$, $(4,4,2)$, $(6,2,2)$ and $(4,2,2,2)$ the zero with limit $-\frac{D+1}{D-1}$ (and its reciprocal) moves away from its limit at the step from genus $0$ to genus~$1$. For $(2,2,2,2)$ and $(2,2,2,2,2)$ it moves away at the steps from genus $0$ to $1$ and from $1$ to $2$.
\end{enumerate}
In particular every zero of every profile of degree at most $10$ moves strictly towards its limit from genus $2$ onwards, and from genus $0$ in odd degree.
\end{theorem}

The exceptional profiles all have only even parts and lie in even degree, and the exceptional zero is always the one whose limit is closest to $-1$. For $(2,2)$ the step from genus $0$ to $1$ starts exactly at the limit.

\begin{proof}
The zeros of $h_r$ come in reciprocal pairs, since $h_r(t)=t^dh_r(1/t)$, and $t\mapsto1/t$ exchanges ``from below towards $x<-1$'' with ``from above towards $1/x\in(-1,0)$''. So it suffices to study the zeros below $-1$. Write $h_r=t(1+t)^{\delta}p_g(t)$ with $\delta=1$ if $D$ is even and $0$ otherwise. Then $p_g$ is palindromic of even degree $2e$, and $e$ is the number of limits below $-1$.

\emph{Exact closed forms.} By Lemma~\ref{lem:poles}, each coefficient of $h_r$ along $r=r_0+2g$, $r_0=d+n-2$, equals $\sum_{k=1}^{D}P_k(g)(k^2)^g$ with $\deg P_k<\lfloor d/(k+1)\rfloor$ for every $g\ge0$. We compute $h_r$ exactly at the first genera from CDM's cut-and-join recursion \cite[Theorem~3.7]{cdm}, solve for the $P_k$ in exact rational arithmetic, and check the result at two further genera.

\emph{Real-rootedness.} Choose rational separators $s_1<\dots<s_{e}=-1$ with $s_i$ halfway between consecutive limits below $-1$. For fixed $s$, $g\mapsto h_r(s)$ is again an exponential polynomial in $g$, and so is $g\mapsto h_r'(-1)$, which replaces $h_r(-1)$ when $\delta=1$. Its dominant term $(D^2)^g$ has a constant nonzero coefficient, proportional to $A_D(s)$ or $A_D'(-1)$. We certify its sign for all $g\ge G$ by an explicit tail bound. Write the exponential polynomial as $\sum u\,g^j\omega^g$ and let $\Omega$ be its largest base, with coefficient polynomial $\sum_jc_jg^j$ of degree $j_0$ (in fact $j_0=0$ in every case). If $c_{j_0}-\sum_{j<j_0}|c_j|/G>0$, if $(1+1/G)^j\omega\le\Omega$ for every other term, and if $(c_{j_0}-\sum_{j<j_0}|c_j|/G)\,G^{j_0}\Omega^G>\sum_{\omega<\Omega}|u|\,G^{j}\omega^{G}$, then the sign holds for every $g\ge G$, because each $g^j(\omega/\Omega)^g$ decreases from $G$ on. The finitely many genera below $G$ are checked exactly. The certified signs alternate, starting from the positive leading coefficient of $p_g$. So for all $g\ge G_0$ (with $G_0\le1$ in every case) $p_g$ has a zero in each window between consecutive separators below $-1$, hence by palindromy $2e$ distinct real zeros. At the genera below $G_0$ the number of distinct real zeros is checked exactly by Sturm sequences.

\emph{Monotonicity on a tail.} Let $x<-1$ be a limit, $b=(D+f-1)(1+x)R_j<0$ the coefficient in \eqref{eq:shift}, $\lambda=q^2$ and $\varepsilon=(1-\lambda)/(3(1+\lambda))$. At the moving points $u_\pm(g)=x+b(1\pm\varepsilon)q^{r_0}\lambda^g$, the value $h_r(u_\pm(g))$ is an exponential polynomial in $g$ with rational bases. The terms with base $D^2$ cancel exactly because $A_D(x)=0$, so the dominant base is $D^2\lambda=(D-1)^2$; by \eqref{eq:single} its coefficient has the sign of $\pm\varepsilon A_D'(x)b$, and we certify these signs for all $g\ge G_1$ as above. So for $g\ge G_1$ the zero in the window of $x$ lies between $u_-(g)$ and $u_+(g)$, these intervals lie in the window, and $\lambda(1+\varepsilon)<1-\varepsilon$ makes the intervals at consecutive genera disjoint and ordered towards $x$. The zero therefore increases strictly to $x$ for $g\ge G_1$. We found $G_1\le17$ in every case.

\emph{The finite part.} For $g\le G_1$ all zeros below $-1$ are isolated in disjoint rational intervals by exact real-root isolation, and consecutive genera are compared. This produces exactly the exceptions listed.
\end{proof}

The census uses only exact integer and rational arithmetic, and it runs in about a minute (Appendix~\ref{app:code}). For degrees $4$ and $5$ an independent certificate, which bounds the cross-differences of coefficients directly, gives the same conclusions.

\begin{conjecture}\label{q:genus2}
Does every zero of every single-profile polynomial move strictly towards its limit from genus $2$ onwards? Does every profile with a zero that moves away have only even parts, including a part equal to $2$?
\end{conjecture}

The converse of the second part fails: $(4)$, $(4,2)$, $(4,4)$ and $(6,2)$ have only even parts and no exception. In the range computed the exceptions are exactly $(2,2)$ and the profiles with only even parts and at least three parts.

\section{Higher terms and short cycles}\label{sec:higher}
This section shows that the early terms of the expansion see only short cycles, and computes the third term for a single profile.

\subsection{Short cycles determine the early terms}
\begin{theorem}[Cycle counts]\label{th:cycles}
Let $d\ge3$ and $0\le k<d/2$. On admissible lengths,
\[
 \frac{d!}{2}\,T_r(t)=\sum_{j=0}^{k}(D-j)^rC_j(t)+O_{d,\boldsymbol\mu}\bigl((1+r)^d(D-k-1)^r\bigr),
\]
where:
\begin{enumerate}
\item each $C_j(t)$ depends only on $d$, $N$ and the counts $m_{i,a}$ with $a\le j$, and for fixed $d,N$ it is a polynomial with rational coefficients in these counts and in $t$;
\item $C_j$ is divisible by $A_{D-j}$, with quotient of degree at most $j$; in particular $C_0=A_D$ and $C_1=A_{D-1}B$;
\item for one profile, with $m_a=m_a(\mu)$, the coefficient of $t$ is
\begin{equation}\label{eq:tangent}
 [t]\,C_j(t)=[y^j]\,(1-y)^{d-2}\prod_{a=1}^{d}(1-y^a)^{m_a},
\end{equation}
which equals $-m_j$ plus a polynomial in $m_1,\dots,m_{j-1}$ when $j\ge1$.
\end{enumerate}
\end{theorem}
\begin{proof}
Fix $j\le k$ and put $\gamma=D-j$. By Lemma~\ref{lem:poles}, a pole at $1/\gamma$ arises only in terms of \eqref{eq:mobius} with a unique block of size at least $\gamma+1=d-j>d/2$, and it is simple. The other blocks contain $a\le j$ letters in total. In the large block a box of content $\gamma$ lies in the first row, so the large diagram has the form $(s-p,\nu)$ with $\nu\vdash p\le j-a$, where $s=d-a$ is the size of the large block.

We use the classical fact that the character of $(s-p,\nu)$ with $|\nu|=p$ is determined by the numbers of cycles of length at most $p$, as a polynomial in these numbers. To see it, expand the Jacobi--Trudi determinant of $s_{(s-p,\nu)}$ into products of complete symmetric functions. In each nonzero term the factor from the first row has index at least $s-p$, and the remaining indices sum to at most $p$. Under the Frobenius characteristic map such a product is the character of a permutation representation on ordered set partitions whose blocks other than the first have at most $p$ letters in total. A permutation fixes such an ordered set partition only if each small block is a union of its cycles. So the character counts choices of short cycles and is a polynomial in their numbers.

Apply this to the large block, with $p\le j-a$. The restrictions to the small blocks use cycles of total size $a\le j$. The probability of a given restriction is a rational constant times a product of falling factorials of the counts $m_{i,b}$ with $b\le a$, and the counts left for the large block are the original counts minus those used. The small-block factors depend only on the small cycle types. So every term contributing at $1/\gamma$, including its Möbius coefficient, is a polynomial in the counts of cycles of length at most $j$. This proves (i). Every contributing large diagram contains the contents $0,1,\dots,\gamma$, so its residue has the factor $A_\gamma(t)$; the small-block factors are polynomials in $t$ at $z=1/\gamma$. This gives divisibility, and the degree bound follows because $T_r$ has degree at most $D$ in $t$. The identifications of $C_0$ and $C_1$ are Theorem~\ref{th:two}. Subtracting the principal parts at these poles and their negatives leaves poles of modulus at least $1/(D-k-1)$, which gives the remainder.

For (iii), the coefficient of $t$ in $h_r$ counts transitive words with one larger letter. Transitivity forces that letter to be $d$, so these are the transitive star factorisations of $\sigma$. In the form of \cite[Eq.~(20)]{clr}, their number is
\[
 \frac{r!}{d!}\,[v^r]\bigl(e^{v/2}-e^{-v/2}\bigr)^{d-2}\prod_{a}\bigl(e^{av/2}-e^{-av/2}\bigr)^{m_a}.
\]
The product equals $e^{Dv}\sum_{i=0}^{2D}c_ie^{-iv}$ with $c_i=[y^i](1-y)^{d-2}\prod_a(1-y^a)^{m_a}$, so $d!\,[t]h_r=\sum_ic_i(D-i)^r$. The polynomial $\sum_ic_iy^i$ satisfies $c_{2D-i}=(-1)^{d+n}c_i$, and $(-1)^{d+n}=(-1)^r$ on admissible lengths, so the bases $D-i$ and $i-D$ contribute equally and $\frac{d!}2[t]h_r=\sum_{i<D}c_i(D-i)^r$ for $r\ge1$. Distinct exponentials are linearly independent, so comparing with the expansion gives \eqref{eq:tangent}. Only the factor $(1-y^j)^{m_j}$ contributes a term $-m_jy^j$ that involves $m_j$.
\end{proof}

For completed-cycle Hurwitz numbers without transitivity, ALR express all coefficients as polynomials in the part multiplicities \cite[Theorem~B]{alr}; part (i) is the analogous statement for the early terms in the connected, $t$-deformed monotone case.

Profiles that agree in their first $k$ cycle counts therefore share their first $k+1$ terms at fixed length. Conversely, by \eqref{eq:tangent} the terms $C_0,\dots,C_s$ with $s=\lfloor(d-1)/2\rfloor$ determine $m_1,\dots,m_s$ one at a time, and with the parity of $\sigma$ this determines $\mu$: for odd $d$ at most one cycle is longer than $s$, and for even $d=2s+2$ the only alternatives with no short cycles are $(d)$ and $(s+1,s+1)$, which have opposite signs.

\subsection{The third term}
\label{sec:third}
This subsection makes Theorem~\ref{th:cycles} explicit for $j=2$ and one profile. Here $f=m_1(\mu)$ and $c=m_2(\mu)$ count fixed points and $2$-cycles, and the third term involves exactly these two numbers.
\begin{theorem}\label{th:third}
For $d\ge5$, the third coefficient in Theorem~\ref{th:cycles} is
\begin{equation}\label{eq:third}
 C_2(t)=A_{d-3}(t)\bigl[a(1+t)^2-Qt\bigr],
\end{equation}
where
\begin{align}
 a&=\binom{d+f-2}{2}-c,\label{eq:third-a}\\
 Q&=\frac{(d-1)\bigl[2c(d-2)+(d-3)(d-4)+2f(d+f-5)\bigr]}
 {(d-4)(d-2)}.\label{eq:third-Q}
\end{align}
Thus, on the admissible parity,
\begin{align*}
 \frac{d!}{2}h_r(\mu;t)
 ={}&(d-1)^rA_{d-1}(t)
 -(d+f-2)(1+t)(d-2)^rA_{d-2}(t)\\
 &+(d-3)^rA_{d-3}(t)\bigl[a(1+t)^2-Qt\bigr]
 +O_{d,\mu}\bigl((1+r)^d(d-4)^r\bigr).
\end{align*}
\end{theorem}
\begin{proof}
Put $\ell=d-3$ and $D=d-1$. At the pole $z=1/\ell$ abbreviate
\[
 F_b=\frac{1+(t-1)b/\ell}{1-b/\ell}.
\]
The generating function for connected words is given by
Lemma~\ref{lem:mobius} and the content product \eqref{eq:char}. A block supporting the pole has size at least $d-2$.
Since $d\ge5$, there is a unique such block, and the pole is simple.
In each contributing diagram the contents $0,1,\ldots,\ell$ supply
the common polar coefficient $A_\ell(t)$. We compute the remaining
multiplier after multiplying the residue by $d!$.

For the one-block partition, the four possible shapes and their
contributions are
\begin{align*}
 (d)       &\quad F_DF_{D-1},\\
 (d-1,1)   &\quad D(f-1)F_{D-1}F_{-1},\\
 (d-2,2)   &\quad \frac{d(d-3)}{2}
                   \left(\binom f2+c-f\right)F_{-1},\\
 (d-2,1,1) &\quad \frac{(d-1)(d-2)}{2}
                   \left(\binom{f-1}{2}-c\right)F_{-1}F_{-2}.
\end{align*}
The character identities in the last two rows follow directly
from Jacobi--Trudi. The identity
$s_{(d-2,2)}=h_{d-2}h_2-h_{d-1}h_1$
counts invariant two-element sets and subtracts fixed points.
The expansion
\[
 s_{(d-2,1,1)}
 =h_{d-2}(h_1^2-h_2)-h_{d-1}h_1+h_d
\]
has character $f(f-1)-\binom f2-c-f+1
=\binom{f-1}{2}-c$.
Here $\binom x2=x(x-1)/2$, including $x=-1$ when $f=0$.
Their values at the identity give the dimensions shown above.

For a singleton and its complement, the singleton must be one of
the $f$ fixed points. The Möbius coefficient is $-1$. The
complementary diagrams with content $\ell$ are $(d-1)$ and
$(d-2,1)$, and the restricted permutation has $f-1$ fixed points.
After the factor $d!/(d-1)!=d$, the contribution is
\[
 -df\bigl[F_{D-1}+(d-2)(f-2)F_{-1}\bigr].
\]
When $f=0$ this term is zero.

The remaining possibilities have two letters outside the large
block. A two-element block is either a pair of fixed points, of
which there are $\binom f2$, or the support of a $2$-cycle, of
which there are $c$. Its unconstrained word generating function
is $(F_1+F_{-1})/2$ in the first case and
$(F_1-F_{-1})/2$ in the second. Both partitions have Möbius
coefficient $-1$. With $u=f(f-1)$ and
$d!/(d-2)!=dD$, their contribution is
\[
 -\frac{dD}{4}\bigl[(u+2c)F_1+(u-2c)F_{-1}\bigr].
\]
Finally, two singleton blocks require two fixed points.
There are $\binom f2$ such partitions, each with Möbius coefficient
$+2$, giving $dD u$. This exhausts all set partitions capable
of supporting the pole.

Let $W(t)$ be the sum of the seven contributions just listed.
The values of the remaining linear factors are
\[
 F_D=1-\frac{d-1}{2}t,\quad F_{D-1}=1-(d-2)t,\quad
 F_{-1}=1-\frac{t}{d-2},\quad
 F_{-2}=1-\frac{2t}{d-1},\quad F_1=1+\frac{t}{d-4}.
\]
Their substitution gives
$[t^0]W=[t^2]W=a$ and $[t]W=2a-Q$,
with $a,Q$ as in \eqref{eq:third-a}--\eqref{eq:third-Q}.
Consequently $C_2=A_\ell W$ has the stated form.
The negative pole contributes equally on admissible lengths,
and all remaining poles have reciprocal magnitudes at most
$d-4$. Partial fractions give the asserted remainder.
\end{proof}

For $d=5$, the full-cycle profile has $a=3$ and $Q=8/3$, giving
$C_2(t)=t(1+t)\bigl[3(1+t)^2-\frac83t\bigr]$.
For $d=6$ and $\mu=(2,2,2)$, the values are $a=3$ and $Q=75/4$.
The requirement $d\ge5$ ensures that $F_1$ is finite at the
third pole and that the large block is unique; for $d=4$ the third pole is $z=1$, where $F_1$ itself has a pole.


\section{A probabilistic reading}\label{sec:prob}
This section interprets the results as statements about a random factorisation. Fix an admissible length $r$, that is, one with $(-1)^r=\sgn\sigma$, and choose a transitive monotone factorisation of $\sigma$ of length $r$ uniformly at random.

\begin{proposition}[Hive number]\label{pr:hive}
As $g\to\infty$, the hive number converges in law to $\sum_{j=1}^{D}\xi_j$ with independent $\xi_j\sim\mathrm{Bernoulli}(j/D)$. Its limiting mean and variance are $(D+1)/2$ and $(D^2-1)/(6D)$. Whenever $h_r(\mu;t)$ is real-rooted, with zeros $0$ and $\alpha_1,\dots,\alpha_{D-1}<0$, the hive number at that length is itself exactly a sum of independent Bernoulli variables with success probabilities $1$ and $p_j=1/(1-\alpha_j)$.
\end{proposition}
\begin{proof}
By \eqref{eq:single} the probability generating function $h_r(t)/h_r(1)$ converges to $A_D(t)/A_D(1)=\prod_{j=1}^D(1-\frac jD+\frac jDt)$. For the second claim factor $h_r(t)/h_r(1)=t\prod_j(t-\alpha_j)/(1-\alpha_j)$ and note $(t-\alpha)/(1-\alpha)=1-p+pt$ with $p=1/(1-\alpha)\in(0,1)$.
\end{proof}

So Theorem~\ref{th:census} says that for every profile of degree at most $10$ the hive number is, at every genus, a sum of independent Bernoulli variables, and Theorem~\ref{th:roots} says that their success probabilities $p_j=1/(1-\alpha_j)$ converge to $j/D$, eventually monotonically, from below for $j<D/2$ and from above for $j>D/2$.

The next result refines the hive-number law by recording which larger letters are used.

\begin{proposition}[Which larger letters are used]\label{pr:occupancy}
Fix $d\ge3$ and $\sigma\in S_d$. For a uniformly random transitive monotone factorisation of $\sigma$ of admissible length $r$, let $\xi_b$ indicate whether $b$ occurs as a larger letter ($2\le b\le d$). As $r\to\infty$ the vector $(\xi_2,\dots,\xi_d)$ converges in law to independent Bernoulli variables with $\Pr(\xi_b=1)=(b-1)/(d-1)$. The largest letter $d$ is always used, since transitivity requires it.
\end{proposition}
\begin{proof}
Give each letter its own marker: weight a word by $\prod_{b\text{ used}}t_b$. As in Lemma~\ref{lem:content} the unconstrained generating function for endpoint $\sigma$ is the coefficient of $\sigma$ in $X(z)=\prod_{b=2}^dF_{J_b}(z;t_b)$, where $F_w(z;s)=(1+(s-1)wz)/(1-wz)$. Now $X$ is not central, so we use Young's orthogonal basis $\{v_T\}$ of each $V_\lambda$, indexed by standard tableaux $T$ of shape $\lambda$, on which $J_bv_T=c_T(b)v_T$, where $c_T(b)$ is the content of the box of $T$ containing $b$. Thus $X$ is diagonal: $Xv_T=\prod_bF_{c_T(b)}(z;t_b)v_T$. For any element $Y$ of the group algebra, $[\sigma]Y=\frac1{d!}\sum_\lambda\dim\lambda\;\mathrm{tr}_{V_\lambda}(\rho_\lambda(\sigma^{-1})Y)$, so
\[
 [\sigma]X=\frac1{d!}\sum_{\lambda\vdash d}\dim\lambda\sum_T\rho_\lambda(\sigma^{-1})_{TT}\prod_{b=2}^dF_{c_T(b)}(z;t_b).
\]
The inclusion--exclusion \eqref{eq:mobius} and the merge argument of Lemma~\ref{lem:mobius} apply to a fixed $\sigma$ (the conditional-uniformity step is not needed); each block factor is the coefficient of the restriction $\sigma|_P$ in the corresponding product for $S_P$, with the markers of the original labels. The pole orders and locations are as in Lemma~\ref{lem:poles}, because they depend only on contents. Content $D$ occurs only for the one-row tableau, where letter $b$ has content $b-1$ and $\rho_{(d)}=1$. Proper set partitions have blocks of at most $d-1$ letters and so contents at most $D-1$. Hence the coefficient of $(1-Dz)^{-1}$ is
\[
 \frac1{d!}\,t_d\prod_{b=2}^{d-1}\frac{(D-b+1)+(b-1)t_b}{D-b+1},
\]
the pole at $-1/D$ (from the one-column tableau, where $\rho=\sgn$) doubles it on admissible lengths, and all other poles have modulus at least $1/(D-1)$. Dividing by the value at $t_2=\dots=t_d=1$ leaves $t_d\prod_{b=2}^{d-1}\bigl(1-\frac{b-1}{D}+\frac{b-1}{D}t_b\bigr)$, the joint generating function of the stated independent Bernoulli variables.
\end{proof}

Setting all $t_b=t$ recovers Proposition~\ref{pr:hive}. The limit does not depend on $\sigma$, although the finite-length occupancy counts do: they are not class functions.

\section{What remains open}\label{sec:open}
This section lists what the results do not settle.

All results here are for fixed degree and growing genus. CDM's Conjectures~3.13 and~3.14 at \emph{every} genus remain open beyond the census of Theorem~\ref{th:census}, which also does not certify interlacing at small genus. The uniform threshold of Theorem~\ref{th:uniform} is far from the truth suggested by the census; a threshold of order $d$ or $\log d$ in the genus, or a proof of Question~\ref{q:genus2}, would be a substantial improvement. The joint regime in which degree and genus grow together is not addressed; for one-part monotone Hurwitz numbers Barazer and Louf \cite{bl} obtain such asymptotics. Coulter and Do \cite{cd} restate the real-rootedness and interlacing conjectures for a $b$-deformation coming from real Grassmannians \cite[Conjectures~4.21 and~4.22]{cd}; the method of Section~\ref{sec:two} should apply there, but we have not carried it out. In the multi-profile case the direction of approach where $B(x_j)=0$ (Remark~\ref{rem:double}) needs the multi-profile analogue of Theorem~\ref{th:third}. Finally, the rate at which the occupancy indicators of Proposition~\ref{pr:occupancy} become independent is given by the next pole, which we have not computed.

\appendix
\section{Reproducible computations}\label{app:code}
This appendix lists the scripts that accompany the paper and what each checks. All arithmetic is exact unless stated. Every check raises an error when it fails, also under \texttt{python -O}. The scripts need Python~3.10 or later, SymPy and mpmath (and NumPy for \texttt{check\_hurwitz.py}).
\begin{itemize}
\item \texttt{verify.py}: direct enumeration of monotone words by increasing larger letter, recording the product and the orbit partition; checks both terms of \eqref{eq:single} for every profile of degrees $2$--$5$ up to length $24$, both terms of \eqref{eq:general} for all unordered pairs of profiles in degrees $3$ and $4$, CDM's cut-and-join recursion against the direct counts, the $(2,2)$ and $(2,2,2)$ polynomials of Theorem~\ref{th:reverse}, the independent all-genus certificates in degrees $4$ and $5$, and the short-cycle dependence of Theorem~\ref{th:cycles} in degrees $6$--$8$.
\item \texttt{check\_third\_residue.py}: the symbolic identity behind Theorem~\ref{th:third} in general $d,f,c,t$, and the theorem for all $55$ profiles of degrees $5$--$8$.
\item \texttt{verify\_hurwitz.py}, \texttt{verify\_multiple.py}: an independent standard-library check of both terms for all profiles of degrees $4$ and $5$ at several values of $t$ ($91$ cases) and of the multi-profile formula for pairs and triples in degree $4$ ($72$ cases).
\item \texttt{check\_hurwitz.py}: an independent permutation-basis dynamic program; filtered two-term identities for $(2,1,1)$ and $(4,2,1)$, and the genus-$20$ zeros of $(4,2,1)$ (CDM's Figure~3).
\item \texttt{figure2.py}: Proposition~\ref{pr:fixed} against direct enumeration ($212$ cases in degrees $3$--$6$), and Table~\ref{tab:fig2}.
\item \texttt{occupancy.py}: Proposition~\ref{pr:occupancy} for every endpoint permutation in degrees $3$--$5$ ($2136$ cases), by direct enumeration and exact exponential-polynomial fitting justified by Lemma~\ref{lem:poles}.
\item \texttt{census.py}: Theorem~\ref{th:census}.
\end{itemize}

\section*{Statements}
\textbf{AI use.} This paper was prepared with substantial assistance from AI systems, which drafted text, proofs and code. The computations were rerun independently from the released scripts. \textbf{Licence.} Text CC0; code MIT.

\begin{thebibliography}{9}
\bibitem{alr} D. Accadia, D. Lewa\'nski and G. Ruzza. On the large genus of Hurwitz numbers. arXiv:2604.12880, 2026.
\bibitem{bl} S. Barazer and B. Louf. High genus one part monotone Hurwitz numbers. arXiv:2604.01190, 2026.
\bibitem{clr} J. Campion Loth and A. Rattan. Centrality of star and monotone factorisations. \textit{Bull. London Math. Soc.} 57(11) (2025) 3567--3585. \url{https://doi.org/10.1112/blms.70170}. Equation numbers refer to arXiv:2403.08354v2.
\bibitem{cd} X. Coulter and N. Do. From Weingarten calculus for real Grassmannians to deformations of monotone Hurwitz numbers and Jucys--Murphy elements. arXiv:2506.04002, 2025.
\bibitem{cdm} X. Coulter, N. Do and E. Moskovsky. Integration on complex Grassmannians, deformed monotone Hurwitz numbers, and interlacing phenomena. \textit{Canad. J. Math.}, published online 2025. \url{https://doi.org/10.4153/S0008414X25101168}.
\bibitem{dhr} N. Do, J. He and H. Robertson. The structure of Hurwitz numbers with fixed ramification profile and varying genus. arXiv:2409.06655, 2024.
\bibitem{yang} C. Yang. The structures of simple Hurwitz numbers and monotone Hurwitz numbers with varying genus. arXiv:2503.01920, 2025.
\end{thebibliography}
\end{document}
